\Legend
小明正在修理一个电子盒。盒中有 $n$ 个从左到右编号为 $1,\ldots,n$ 的元件，有 $m$ 条导线连接不同元件。导线 $(x_i,y_i)$ 满足 $x_i<y_i$，且同一对元件之间至多有一条导线。导线无方向，并且任意两条导线不相交；允许在公共端点相接。形式化地，不存在 $i<j<k<l$ 使 $(i,k)$ 与 $(j,l)$ 同时存在。

检测仪有四个探针。依次接到四个两两不同的元件 $i,j,k,l$ 后，若 $(i,j)$、$(j,k)$、$(k,l)$ 均有导线，则检测仪可以正常工作。

接下来会逐一加入 $q$ 条导线 $(u_i,v_i)$，保证加入后仍满足以上全部要求。请输出初始状态以及每次加线后的合法\textbf{有序}四元组数量。

\Input

第一行两个整数 $n,m$，分别表示元件和初始导线的数量。

接下来 $m$ 行，每行两个整数 $x_i,y_i$，表示第 $i$ 条初始导线。

下一行一个整数 $q$，表示加线次数。

接下来 $q$ 行，每行两个整数 $u_i,v_i$，表示第 $i$ 次加入的导线。

\Output

共输出 $q+1$ 行：第一行是初始答案，第 $i+1$ 行是加入第 $i$ 条导线后的答案。

\textbf{数据范围}

$2\le n\le5\times10^5$，$0\le m,q\le10^6$；所有端点满足 $1\le x_i<y_i\le n$、$1\le u_i<v_i\le n$，且导线两两不同、始终无交叉。时间限制 4 秒，空间限制 512 MiB。

\Examples

\begin{example}
\exmp{4 3
1 2
2 3
3 4
2
2 4
1 4}{2
4
12}
\end{example}

最初只有 $(1,2,3,4)$ 及其逆序。第一次加线后新增两个合法四元组；第二次加线后共有十二个。
